\section{September 1, 2026}

Recall that a measure on a $\sigma$-algebra $\mathcal M$ over a set $X$
is a map
\[
  \mu\colon\mathcal M\longrightarrow[0,\infty]
\]
such that $\mu(\varnothing)=0$ and
\[
  \mu\left(\bigcup_{i=1}^{\infty}A_i\right)
  =\sum_{i=1}^{\infty}\mu(A_i)
\]
whenever $A_1,A_2,\ldots\in\mathcal M$ are pairwise disjoint. For
Lebesgue measure $m$ on $\bb{R}$, the normalization is $m([0,1])=1$.

\subsection{Complete measure spaces}

\begin{definition}[Complete measure]
  A measure $\mu$ on a $\sigma$-algebra $\mathcal M$ over $X$ is
  \emph{complete} if every subset of a null set is measurable. That is,
  if $E\subseteq N$ for some $N\in\mathcal M$ with $\mu(N)=0$, then
  $E\in\mathcal M$.
\end{definition}

By monotonicity, every such set $E$ also satisfies
\[
  0\leq\mu(E)\leq\mu(N)=0,
\]
so $\mu(E)=0$.

Every measure space has a canonical completion obtained by adjoining all
subsets of its null sets.

\begin{proposition}[Completion of a measure]
  Let $(X,\mathcal M,\mu)$ be a measure space, and let
  \[
    \mathcal N=\{N\in\mathcal M:\mu(N)=0\}.
  \]
  Define
  \[
    \overline{\mathcal M}
    =\{E\cup F:E\in\mathcal M,\ F\subseteq N
      \text{ for some }N\in\mathcal N\}.
  \]
  Then the following statements hold.
  \begin{enumerate}[label=(\roman*)]
    \item The collection $\overline{\mathcal M}$ is a $\sigma$-algebra
      on $X$ containing $\mathcal M$.
    \item There is a unique measure $\overline\mu$ on
      $\overline{\mathcal M}$ extending $\mu$. It is given by
      \[
        \overline\mu(E\cup F)=\mu(E),
      \]
      where $E\in\mathcal M$ and $F$ is contained in a member of
      $\mathcal N$. Moreover, $\overline\mu$ is complete.
  \end{enumerate}
\end{proposition}

\begin{proof}
  The collections $\mathcal M$ and $\mathcal N$ are closed under
  countable unions. For $\mathcal N$, this follows from countable
  subadditivity: if $N_i\in\mathcal N$, then
  \[
    0\leq\mu\left(\bigcup_{i=1}^{\infty}N_i\right)
    \leq\sum_{i=1}^{\infty}\mu(N_i)=0.
  \]
  Hence $\overline{\mathcal M}$ is closed under countable unions. Indeed,
  if $E_i\in\mathcal M$ and $F_i\subseteq N_i\in\mathcal N$, then
  \[
    \bigcup_{i=1}^{\infty}(E_i\cup F_i)
    =\left(\bigcup_{i=1}^{\infty}E_i\right)
      \cup\left(\bigcup_{i=1}^{\infty}F_i\right),
  \]
  and $\bigcup_iF_i\subseteq\bigcup_iN_i\in\mathcal N$.

  To check closure under complements, let $E\in\mathcal M$ and
  $F\subseteq N\in\mathcal N$. Replacing $N$ by $N\setminus E$ and
  $F$ by $F\setminus E$ does not change $E\cup F$, so we may assume that
  $E\cap N=\varnothing$. In that case,
  \[
    (E\cup F)^c=(E\cup N)^c\cup(N\setminus F).
  \]
  The first set on the right belongs to $\mathcal M$, and the second is
  contained in the null set $N$. Thus $(E\cup F)^c$ belongs to
  $\overline{\mathcal M}$. Since $\varnothing\in\overline{\mathcal M}$,
  this proves that $\overline{\mathcal M}$ is a $\sigma$-algebra. Taking
  $F=\varnothing$ shows that $\mathcal M\subseteq\overline{\mathcal M}$.

  We next show that the formula for $\overline\mu$ is well defined.
  Suppose
  \[
    E_1\cup F_1=E_2\cup F_2,
    \qquad
    E_i\in\mathcal M,
    \quad F_i\subseteq N_i\in\mathcal N.
  \]
  Then $E_1\subseteq E_2\cup N_2$, and therefore
  \[
    \mu(E_1)
    \leq\mu(E_2\cup N_2)
    \leq\mu(E_2)+\mu(N_2)
    =\mu(E_2).
  \]
  By symmetry, $\mu(E_2)\leq\mu(E_1)$, so
  $\mu(E_1)=\mu(E_2)$.

  Clearly $\overline\mu(\varnothing)=0$. Suppose that
  $A_i=E_i\cup F_i$ are pairwise disjoint members of
  $\overline{\mathcal M}$, where $E_i\in\mathcal M$ and
  $F_i\subseteq N_i\in\mathcal N$. The sets $E_i$ are pairwise disjoint,
  so
  \begin{align*}
    \overline\mu\left(\bigcup_{i=1}^{\infty}A_i\right)
    &=\mu\left(\bigcup_{i=1}^{\infty}E_i\right)\\
    &=\sum_{i=1}^{\infty}\mu(E_i)\\
    &=\sum_{i=1}^{\infty}\overline\mu(A_i).
  \end{align*}
  Thus $\overline\mu$ is a measure, and it extends $\mu$.

  If $A=E\cup F\in\overline{\mathcal M}$ and
  $\overline\mu(A)=0$, then $\mu(E)=0$. Since
  $A\subseteq E\cup N$ and $E\cup N$ is a measurable null set, every
  subset of $A$ belongs to $\overline{\mathcal M}$. Hence
  $\overline\mu$ is complete.

  Finally, let $\nu$ be any measure on $\overline{\mathcal M}$ that
  extends $\mu$. If $A=E\cup F$, with $E\in\mathcal M$ and
  $F\subseteq N\in\mathcal N$, then
  \[
    E\subseteq A\subseteq E\cup N.
  \]
  Monotonicity and the fact that $\nu$ extends $\mu$ give
  \[
    \mu(E)=\nu(E)\leq\nu(A)\leq\nu(E\cup N)
    =\mu(E\cup N)=\mu(E).
  \]
  Thus $\nu(A)=\mu(E)=\overline\mu(A)$, which proves uniqueness.
\end{proof}

The measure $\overline\mu$ is called the \emph{completion} of $\mu$.

\subsection{Outer measures}

The goal is to construct measures from simple geometric pieces. For
example, to estimate the area of a set $E\subseteq\bb{R}^2$, one can
cover $E$ by squares of side length $s$ and consider
\[
  \inf_{s>0}\bigl(\text{number of squares in the cover}\bigr)s^2.
\]
More generally, one covers a set by countably many simple pieces and
takes the infimum of the sums of their assigned sizes.

A related guiding principle is that the ``reasonable'' subsets
$E\subseteq[0,1]^2$ should be those for which outer and inner area agree.
If the unit square has area $1$, its inner estimate for $E$ is
\[
  1-m^*\bigl([0,1]^2\setminus E\bigr).
\]

\begin{definition}[Outer measure]
  An \emph{outer measure} on a set $X$ is a map
  \[
    \mu^*\colon 2^X\longrightarrow[0,\infty]
  \]
  satisfying the following properties.
  \begin{enumerate}[label=(\roman*)]
    \item $\mu^*(\varnothing)=0$.
    \item If $A\subseteq B\subseteq X$, then
      $\mu^*(A)\leq\mu^*(B)$.
    \item For every sequence $A_1,A_2,\ldots\subseteq X$,
      \[
        \mu^*\left(\bigcup_{j=1}^{\infty}A_j\right)
        \leq\sum_{j=1}^{\infty}\mu^*(A_j).
      \]
  \end{enumerate}
\end{definition}

Unlike a measure, an outer measure is defined on every subset of $X$;
its countable subadditivity therefore applies without a measurability
assumption.

\begin{proposition}[Outer measure from coverings]
  Let $\mathcal E\subseteq 2^X$ contain $\varnothing$ and $X$, and let
  \[
    \rho\colon\mathcal E\longrightarrow[0,\infty]
  \]
  satisfy $\rho(\varnothing)=0$. For $A\subseteq X$, define
  \[
    \mu^*(A)
    =\inf\left\{
      \sum_{j=1}^{\infty}\rho(E_j):
      E_j\in\mathcal E,\quad
      A\subseteq\bigcup_{j=1}^{\infty}E_j
    \right\}.
  \]
  Then $\mu^*$ is an outer measure on $X$.
\end{proposition}

\begin{remark}
  Even when $E\in\mathcal E$, it need not be true that
  $\mu^*(E)=\rho(E)$. A more efficient countable cover may have smaller
  total $\rho$-size than the single-set cover by $E$.
\end{remark}

\begin{example}
  Let $X=\bb{R}$ and take
  \[
    \mathcal E
    =\{\varnothing,\bb{R}\}
      \cup\{[a,b]:a,b\in\bb{R},\ a\leq b\}.
  \]
  Define
  \[
    \rho(\varnothing)=0,
    \qquad
    \rho(\bb{R})=\infty,
    \qquad
    \rho([a,b])=(b-a)^2.
  \]
  Although $\rho([0,1])=1$, the induced outer measure satisfies
  $\mu^*([0,1])=0$. Indeed, for every $n\in\bb{N}$,
  \[
    [0,1]
    =\bigcup_{j=0}^{n-1}
      \left[\frac{j}{n},\frac{j+1}{n}\right],
  \]
  and the total $\rho$-size of this cover is
  \[
    \sum_{j=0}^{n-1}\left(\frac1n\right)^2
    =\frac1n.
  \]
  Hence $0\leq\mu^*([0,1])\leq 1/n$ for every $n$, which gives the
  claim.
\end{example}

\begin{proof}[Proof of the proposition]
  The defining collection of sums is nonempty because any set
  $A\subseteq X$ is covered by the sequence $X,\varnothing,
  \varnothing,\ldots$. Thus $\mu^*(A)$ is well defined as an element of
  $[0,\infty]$.

  The sequence $\varnothing,\varnothing,\ldots$ covers $\varnothing$,
  so
  \[
    0\leq\mu^*(\varnothing)
    \leq\sum_{j=1}^{\infty}\rho(\varnothing)=0.
  \]

  If $A\subseteq B$, then every countable cover of $B$ by members of
  $\mathcal E$ is also a cover of $A$. Taking infima over the respective
  collections of covers gives $\mu^*(A)\leq\mu^*(B)$.

  It remains to prove countable subadditivity. Let $A_j\subseteq X$ for
  $j\in\bb{N}$. If $\sum_j\mu^*(A_j)=\infty$, there is nothing to prove.
  Otherwise, fix $\varepsilon>0$. For every $j$, choose
  $E_{j,k}\in\mathcal E$, $k\in\bb{N}$, such that
  \[
    A_j\subseteq\bigcup_{k=1}^{\infty}E_{j,k}
    \qquad\text{and}\qquad
    \sum_{k=1}^{\infty}\rho(E_{j,k})
    \leq\mu^*(A_j)+\varepsilon 2^{-j}.
  \]
  The doubly indexed collection $\{E_{j,k}:j,k\in\bb{N}\}$ is
  countable and covers $\bigcup_jA_j$. Consequently,
  \begin{align*}
    \mu^*\left(\bigcup_{j=1}^{\infty}A_j\right)
    &\leq\sum_{j=1}^{\infty}\sum_{k=1}^{\infty}\rho(E_{j,k})\\
    &\leq\sum_{j=1}^{\infty}
      \left(\mu^*(A_j)+\varepsilon 2^{-j}\right)\\
    &=\sum_{j=1}^{\infty}\mu^*(A_j)+\varepsilon.
  \end{align*}
  Letting $\varepsilon\downarrow0$ proves the result.
\end{proof}

\begin{definition}[Measurability with respect to an outer measure]
  Let $\mu^*$ be an outer measure on $X$. A set $A\subseteq X$ is
  \emph{$\mu^*$-measurable} if, for every $E\subseteq X$,
  \[
    \mu^*(E)
    =\mu^*(E\cap A)+\mu^*(E\cap A^c).
  \]
\end{definition}

Countable subadditivity always gives
\[
  \mu^*(E)
  \leq\mu^*(E\cap A)+\mu^*(E\cap A^c),
\]
so the content of the measurability criterion is the reverse inequality.
When $\mu^*(E)<\infty$, the quantity
\[
  \mu^*(E)-\mu^*(E\cap A^c)
\]
can be viewed as the inner measure of $E\cap A$ relative to $E$.
